\section{\texorpdfstring{Plovoucí řádová čárka}{Plovouci radova carka}}

% =====================================================================================================================
  \begin{frame}
    \frametitle{Pevná řádová čárka X Plovoucí řádová čárka}
    \boldmath
    
    \begin{itemize}
      \item Pevná řádová čárka má velmi omezený rozsah čísel \\
      \textbf{Příklad (16 bit) $\Longrightarrow$ jen 65536 zobrazitelných hodnot, tedy maximální rozsah hodnot je pro:}
      \vspace{2mm}
      \begin{itemize}
        \item celá čísla bez znaménka: \textcolor{blue}{$0$} až \textcolor{blue}{$65536$},
        \item celá čísla se znaménkem: \textcolor{blue}{$-32768$} až \textcolor{blue}{$32767$},
      \end{itemize}
      Po zavedení neceločíselné části se sníží jak jednotka mřížky, tak maximální rozsah hodnot.
      \vspace{3mm}
      \item Rozsah lze zvětšit vědeckým zápisem čísel (vědecká notace)
      \begin{center}
        \textcolor{blue}{$m\cdot Z^{e}$}
      \end{center}
    \end{itemize}
	\end{frame}

% =====================================================================================================================
  \begin{frame}
    \frametitle{Zápis}
    \boldmath
    
    Ve všech soustavách stejný:
    \large
    \begin{center}
      \framebox{$m\cdot Z^{e}$}
    \end{center}
    
    \normalsize
    \begin{itemize}
      \item $m$ ... \textbf{MANTISA}, informace o číselné hodnotě
      \item $e$ ... \textbf{EXPONENT}, informace o pozici řádové čárky
      \item $Z$ ... \textbf{ZÁKLAD}, základ soustavy vyjadřovaného čísla
    \end{itemize}
    
    \vspace{3mm}
    Mantisa i exponent mohou nabývat kladných i záporných hodnot.
    
    \vspace{2mm}
    \textbf{Příklad v dekadické soustavě pro 2 číslice mantisy a 1 číslici exponentu obojí s přímým zápisem znaménka:}
    
    \vspace{2mm}
    \begin{itemize}
      \item rozsah hodnot: \textcolor{blue}{$-9,9\cdot 10^{9}$} až \textcolor{blue}{$9,9\cdot 10^{9}$}
      \item jednotka, krok: \textcolor{blue}{$0,1\cdot 10^{-9}$}
    \end{itemize}
	\end{frame}
% =====================================================================================================================
  \begin{frame}
    \frametitle{Vlastnosti}
    \boldmath
    \textbf{Vyjádřením čísel v plovoucí řádové čárce}
    \begin{itemize}
      \item se nemění počet zobrazitelných hodnot. Pro 16 bitů bude stále počet zobrazitelných hodnot roven 65536,
      \item lze získat větší rozsah hodnot z pohledu maxima a minima,
      \item je přesnost každého čísla dána jednotkou zápisu mantisy,
      \item je třeba vyřešit problém nejednoznačnosti vyjádření téhož čísla, příklad pro dekadickou hodnotu:
      \small
        \begin{align*}
        62,15 =& 62,15 \cdot 10^{0} = 621,5 \cdot 10^{-1} = 6,215 \cdot 10^{1} = ...
        \end{align*}
      \normalsize
      \item musíme upravit aritmetické operace tak, aby respektovaly přítomnost exponentu, viz dále.
    \end{itemize}
	\end{frame}
  
% =====================================================================================================================
  \begin{frame}
    \frametitle{ANSI/IEEE 754 - kódování hodnoty}
    \boldmath
    \textbf{Základní přesnost 32 bit:}
    \vspace{3mm}
    
    \begin{minipage}[t][0.2\textheight][t]{0.15\textwidth}
    \begin{center}
      \textcolor{blue}{\textbf{1 bit}}
      
      \begin{tabular}{|c|}
      \hline
      s \\
      \hline
      \end{tabular}
      
      znaménko
    \end{center}
    \end{minipage}
    \begin{minipage}[t][0.2\textheight][t]{0.5\textwidth}
    \begin{center}
      \textcolor{blue}{\textbf{8 bit}}
      
      \begin{tabular}{|c|c|c|c|c|c|c|c|}
      \hline
      e & e & e & e & e & e & e & e \\
      \hline
      \end{tabular}
      
      exponent
    \end{center}
    \end{minipage}
    \begin{minipage}[t][0.2\textheight][t]{0.3\textwidth}
    \begin{center}
      \textcolor{blue}{\textbf{23+1 bit}}
      
      \begin{tabular}{|c|c|c|c}
      \hline
      m & m & m & ...\\
      \hline
      \end{tabular}
      
      mantisa
    \end{center}
    \end{minipage}
    
    \begin{itemize}
      \item Znaménko náleží mantise - přímý kód:
      
      \begin{itemize}
        \item 0 ... kladné
        \item 1 ... záporné
      \end{itemize}
      
      \item Exponent je zapisován aditivním kódu s posunutím \textcolor{red}{\textbf{+127}}.
      \item Mantisa je zapisována v normalizovaném tvaru, který obsahuje \textcolor{red}{\textbf{skrytou 1}}. Mantisa je tak o jeden bit delší než kolik je počet bitů vyhrazených pro její zápis.
    \end{itemize}

	\end{frame}
  
% =====================================================================================================================
  \begin{frame}
    \frametitle{Aritmetika}

    \begin{enumerate}
      \item Aritmetická operace
      
      \begin{itemize}
        \item \textbf{sčítání (odečítání)}
        
        \begin{enumerate}
          \item Úprava zápisu obou hodnot tak, aby měly shodné exponenty.
          \item Sečtení (odečtení) mantis v přímém kódu.
        \end{enumerate}
        
        \item \textbf{násobení}
        \begin{enumerate}
          \item Součet exponentů v aditivním kódu.
          \item Násobení mantis.
          \item Násobení znamének - jde vlastně o sečtení znamének, funkce XOR.
        \end{enumerate}
        
        \item \textbf{dělení}
        \begin{enumerate}
          \item Rozdíl exponentů v aditivním kódu.
          \item Dělení mantis.
          \item Totéž co pro násobení znamének.
        \end{enumerate}
        
      \end{itemize}
      
      \item Normalizace výsledku
    \end{enumerate}
	\end{frame}